\section{The Creative Power of the Individual}

Evola concludes his discussion of privation with the observation that it is the Will that throws light on the unknown. The accumulation of knowledge is insufficient.

There is a continuum between spontaneity and free will. Everything objective and necessary in experience --- that is, everything experienced as ``privation'', arises from spontaneity. That means, there is no consciousness of having ``willed'' it; rather, it arises out of the Absolute as the experience of the potential becoming actual. From the perspective of the individual, it arrives ``out of the dark''. Evola elaborates:

\begin{quotex}
Hence the conception that appears in the third stage of the development of the individual is, as a whole, the following: a continuum of activity that has as limits spontaneity on one side and free will on the other. Spontaneity is the universal, free will is the individual. These limits are related to each other as potential to act: everything objective, immediate, and necessary in experience, is, in regards to the position of the individual, the non-being inherent in what is in potential --- and here perhaps we can understand what certain mystics alluded to when they spoke of the ``dark passion of the world'', the ``unspeakable suffering of existence'', in which the body of the ``celestial man'' is crucified.
\end{quotex}


With free will, this darkness is illumined. The ``act'' of the will makes the world ``real'', whereas privation is unreal.

\emph{The individual, now conscious of his power, is creative}.


This brings up three points

\begin{itemize}


\item The individual the source of the becoming of the world, thereby giving it its solidity and reality.

\item This is why the search for certainty in the things of the world will always end in frustration.

\item The individual, at this point, is the source of his own being.
\end{itemize}


He explains:

\begin{quotex}
Freedom is the act and the luminous flame of such a darkness, of such a privation; and the world becomes, it is made real in conformity with absolute reality only in and through this flame --- that is, only to the extent that the individual, asserting himself at the point of power and control, consumes, burns up his former nature that was spontaneously made.
\end{quotex}


Paradoxically, the so-called Traditionalist does not long for a return to the Past; rather it is only in creating the Future does he come to know his own powers and the Will of God Providence . Hence true explanations do not lie in things or in the past, but rather ahead, since they are dependent on the individual for their \emph{raison d'être}. Such an individual is his ``own property'' and ``persuaded/convinced''. (The latter two expressions are due to \textbf{Max Stirner} and \textbf{Carlo Michelstaedter} respectively.)


Evola then claims that it is the individual who gives meaning and purpose to the world. This distinguishes him from Schopenhauer's position that denies any such purpose; in the latter's system, the individual is not free and does not become aware of the noumenal ``Will'' as his own centre. (I think it should be clear by now that Evola is not exalting the ``human being'' as such, since the states he describes are beyond the human state.)

\begin{quotex}
only at the point at which the individual realises himself in the sudden intuition of power does a purpose, a reason and a goal in nature arise: not before; it is he who gives it to them. It demands it of his action. And yet the individual has a single imperative by this point.
\end{quotex}


Evola has been describing the Traditional esoteric path. It consists of three stages.

\begin{description}

\item[Catharsis] This is the stage of purification, moving from spontaneity to Will.
\item[Theoria] This is the stage of illumination, where the darkness is illumined.
\item[Theosis] This is the stage of divinization, where one incarnates God.
\end{description}


To some it may seem that Evola is inverting the process of salvation in the Christian sense. Others may recognize in it the Traditional teachings of the Greek fathers. Evola concludes:

Be, make yourself GOD, and in bringing that about, SAVE the world.

\flrightit{Posted on 2011-03-07 by Cologero}

\begin{center}* * *\end{center}

\begin{footnotesize}
\begin{sffamily}
\texttt{Aidan on 2011-03-08 at 14:07 said:}

Supra-individual states are not associated with individuality. This is not a linguistic turn but a point that is mathematical in its precision. Entry into supra-individual states requires the complete surrender of all individuality -- the partial preservation and integration of specific elements of the latter within the former is a matter which is not in the hands of the individual.

``It is almost superfluous to stress how little place the individual `self' occupies in the totality of the being, since even given its entire extension when envisaged in its integrality, and not merely in one particular modality such as the corporeal, it constitutes only one state like the others, among an indefinitude of others.''

and

``... there is assuredly a form of consciousness, among all possible forms, that is properly human, and this determinate form (ahankara, or `self-consciousness') is inherent in what we term the `mental' faculty, that is, precisely that `internal sense' which is designated in Sanskrit by the name manas, and which is truly the characteristic of human individuality. This faculty ... must be carefully distinguished from pure intellect, which latter, on the contrary, by reason of its universality, must be regarded as existing in all beings and in all states, whatever the modalities through which its existence is manifested,; and one should not see in the `mental' anything beyond what it really is, that is, to use the terminology of logic, a `specific difference' pure and simple, the possession of which does not itself confer on man any effective superiority over other beings.''

Both of these quotes are taken from Guenon's ``The Multiple States of the Being''. If the human state, the individual state and the psychic principle (both mental and sentimental) are all one and the same, how is the individual `asserting himself', `giving meaning and purpose to the world' and `realising himself' not an exaltation of the human state and how are the states Evola describes *beyond* the human state, as opposed to extremities that may be reached within it?

\hfill

\texttt{Cologero on 2011-03-09 at 07:45 said:}

As for ``realising himself'', how is that different from how Guenon describes the ``True Man'', the one who realises all his possibilities? Are you suggesting that the world has no meaning or purpose, that a man acts randomly without meaning, purpose or intent? ``Like, whatever
... '' is not a heroic attitude.


Evola's concern is with the ``I'', which is not individual (``not multipliable''). Compare with Coomaraswamy's essay ``The One and Only Transmigrant''. We've emphasized that there is a difference in perspective between Guenon and Evola. Guenon is strictly a metaphysician, interested in higher states, or really the Transcendent Man, beyond all states. But even the jivanmukti --- the one who has transcended while in this world --- nevertheless continues to act in the world. Is his action random, without purpose, dependent on the external environmental factors of gross and subtle manifestation? Or does he act from the ``center'', the Unmoved Mover? Even the jivanmukti can't avoid ``asserting himself'', just his manner of self-assertion.

Of course, ultimately, there is no ``external environment''. The concept of privation is not a novelty from Evola, even Guenon uses that term.

In Multiple States, Guenon writes about the angelic hierarchy, only to dismiss it. From a metaphysical point of view, all that matters is the existence of multiple states in general, not what they are specifically. However, from a cosmological point of view (and Hermetic, which is a cosmology), these particular states are of great interest.

As for the Individual and the Absolute, Evola writes: ``between the person and universal subject there is a relationship not of coexistence and otherness, rather of continuity and progressivity.'' We will bring that up in the discussion of the One and the Many.

\hfill

\texttt{Aidan on 2011-03-09 at 21:00 said:}

It is frankly astonishing that you can derive any suggestion of an attitude of ``like, whatever'' from my previous comment. Given the reference to individuality being oriented towards and grounded within supra-individual states -- a perspective which transforms both world and individual together -- where on earth do you find even a hint that the world might be without meaning or purpose, or that a man ``acts randomly without meaning, purpose or intent''? When every last part of its order is derived from a still higher order, the world is saturated with meaning and demonstrates perfect purpose; upon the plane of the individual, precisely the same is true of actions which arise directly and spontaneously from the stillness of immovable principle and return to it with complete circularity.

There is no comparison between heroism and divinity. There is not one hero; there are many. There are not many divinities; there is one. Heroism is necessarily attached to individuality while divinity transcends it completely. Compare with Coomaraswamy's `One and Only Transmigrant' indeed! Who is it who transmigrates from incarnation to incarnation? It is not Cologero or Aidan or Evola or Guenon: ``the Lord is the only transmigrant.''

Coomaraswamy and Guenon are at one on this point because both are at one with primordial tradition and their knowledge lacks nothing. Evola, by contrast, is out on his own -- an island in the ocean -- at no point is he talking about the Lord. He means -- as you say -- make YOURSELF God. He means himself, nothing higher: for him, his interior sense is the highest, hence all the sound and fury. It is not a subtle mystery; indeed, it is a crude caricature of true knowledge. Even if he is not oriented towards unity with the Lord, as a brahmin would be (``Whoever is joined unto the Lord is One Spirit'' 1 Corinthians 6:17), he still needs to bow to the Lord, as all kshatriya must, if his deeds are to be genuinely heroic and his heart truly noble.

Yes, the True Man realises himself and all his possibilities but not at the expense of transcendence. He does not have to be transcendent himself -- which is just as well, because he cannot be -- but his relationship to the Transcendent cannot be aberrant if he is to be and remain True.

There would be no problem if there were no insistence that Evola and Guenon are providing related and compatible perspectives. In a fundamental sense, they reveal themselves to be irreconcilable. Either Evola disrupts and deforms Guenon's metaphysics or Guenon severely restricts Evola's ambition. Pressing them together without allowing one to fully dominate the other, as though your own third perspective went further than either, will do a disservice to one or the other and conceal elements in both that need to be seen.

The suggestion that was made earlier that knowledge and action are simply two different spiritual paths is misleading in the extreme. No-one would suggest Evola is a representative of karma marga, a path of selfless action, yet this is the only sense in which action and knowledge can be presented as parallel contrasting paths about which individual choices can be made. Karma marga can be contrasted with jnana marga and bhakti marga without any confusion. The issue of the brahmins and kshatriyas -- and the type of knowledge and type of action they represent -- is not a matter of two different paths that an individual may or may not take. These are social castes, not individual paths. It is the ordering of the world that is at stake. Can there be two world orders simultaneously? This is what Evola disingenuously presents before contradicting himself (and placing action above knowledge) and you echo it. Guenon presents knowledge and action as two aspects of the same order and has knowledge superior to action in the sense that correct and appropriate action proceeds out of knowledge and not vice versa. Both cannot be true; the issue can be fudged, however, which makes it come out in Evola's favour by default. Evola's is a recipe for a kingdom divided against itself and it will not stand -- it will be perpetual all-out war followed by collapse and destruction.

Finally, with regard to the angelic hierarchies, Guenon says that ``the existence of extra-human and supra-human states is of little importance to us'' in the sense that we have ``no motive for occupying ourselves especially with them''. He does not dismiss them at all, speaking of angels as beings in their own right having none of the faculties of an individual order, e.g., reason, but in no way being inferior as a result of this and saying that ``the `angelic' states are properly the supra-individual states of manifestation, that is, those pertaining to non-formal manifestation.''

I intend this to be my last comment on this site. The manner in which my earlier comments have been received has spoken volumes and it would be futile to continue in the same vein. In any event, allow me to thank you for your endeavors. I wish you well.

\hfill

\texttt{Matt on 2011-03-09 at 22:15 said:}

Aidan,

I've felt when reading Evola that he places action above knowledge. He may focus on action a lot more , but I think that isdue to his particular persuasion.

Anyways, hopefully you continue to comment. I think your posts are well-written and contian some good points.

\hfill

\texttt{Matt on 2011-03-09 at 22:16 said:}

Correction: First sentence should be... .I've never*

\hfill

\texttt{Cologero on 2011-03-10 at 08:09 said:}

It is important not to equate ``the courage of one's convictions'' with being possessed by a ``fixed idea''. Although this is a distasteful task, nevertheless I need to clarify some points lest anyone think there is any substance to the objections made to the commentaries on Evola's book. (yes, intelligent objections are welcome.)

The commenter objected to these phrases: `` `asserting himself', `giving meaning and purpose to the world' and `realising himself' ''

Gornahoor intersperses other material with the commentary on Evola for a specific purpose: viz, to situate Evola's ideas within a larger perspective. In that way we can properly understand those phrases, dealing with their actual content rather than nitpicking over ``imperfections of expression.'' Thus:

``Asserting oneself'' needs to be understood as the role of Man as the intermediary between Heaven and Earth; as the Will as intermediary between Providence and Destiny.

If the commenter objects to the perfectly legitimate idea of `giving meaning' how else are we to understand that objection? Is it really preferable that the ``universe'' provides meaning that we passively accept? It is just that attitude that Evola opposes.

``Realising oneself'' is again relatable to Guenon. The individual as Spirit (a non-formal state) actualizes or realizes his potentialities in the World. Again, nothing at all objectionable from a Traditional perspective.

Evola is providing a phenomenological study of the ``I'', that is available to anyone making the effort. He it not writing about the will of Evola, but rather the will of the I. In this he should be understood as taking Shankara's perspective against Ramanuja. It is a legitimate point of view, which Evola develops to its logical conclusion. Again, we pointed out that the concept of privation is common to Aquinas and Guenon. Now, of course, anyone can prefer the position of the ``realist'' to the ``idealist'' (as Evola uses those terms). However, to make the discussion worthwhile, he needs to respond to Evola's arguments against the realist, not to wish that they had never been made.

The idea of the Absolute I, or Brahman, is not unique to Evola; perhaps readers are not so accustomed to dealing with the logical consequences of this view.

It is not so absurd, and subjectivity is not a counter-argument. Go back and read what we have written about the Primordial State, particularly as historically documented by Julian Jaynes. Subjectivity and sense of ego are not necessary or even natural to man. Jaynes documents how at some stage, men were not plagued by such feelings. Instead, they acted together as a unified group. There is a reason for all these posts, they fit together.

Finally, we can mention that there is one Hero, with a thousand faces. Someone once wrote that there are just a few causes of human misery --- illness, poverty, oppression --- and are common to every man. However, the Heroic path is unique to each man. Unlike the path to misery, there is no one way to God, or Enlightenment. Instead each man has to find his own way and deal with his own Destiny. That is why the Hero has a thousand faces.

\hfill

\texttt{Albion on 2011-03-10 at 08:42 said:}

Hi, it's good to be allowed to ask questions, and thank you for this marvelous site.

My question is, if Guenon (I don't have the exact reference) was worried about ``spiritual counterfeits'' and ``amalgams'' that had the form of traditionalism, but whose principle was entirely demonic in its unity, why was it that the ``West'' immediately sprang to mind in the sense that the history of the West for the last thousand years seems to be one entirely of REVOLUTION and re-synthesis which always aspires to conserve? I don't know enough about the East to ask a similar question, but the East seems to me to be lacking in some form of energy. Is this merely my perception as a Westerner under the shadow?

\hfill

\texttt{Graham on 2011-03-10 at 12:04 said:}

I for one cannot understand Aidan's objections. The theory involved is quite easily followed: Evola took the `raw material' of German Idealism and transmuted it into the theosis of a Meister Eckhart or a Plotinus.

\hfill

\texttt{Cologero on 2011-03-10 at 12:09 said:}

I don't understand the point, Albion. You say the west is wracked by REVOLUTION, but where do we see the ``re-synthesis which always aspires to conserve''?

Read the alternatives mentioned by Guenon:
\textit{The Prophecies of Rene Guenon.}\footnote{\url{http://www.gornahoor.net/?p=35}} One of them involves a revolution. Instead of the constant harping about Guenon vs Evola, I wonder if anyone takes any of this seriously? Or is it more fun to re-arrange the deck chairs?

\hfill

\texttt{Graham on 2011-03-10 at 12:30 said:}

Furthermore, if at a *certain stage* in this process one winds up alone like Robinson Crusoe, how can that be construed as an objection to the process itself? Doesn't St. John of the Cross speak of a dark night of the soul? Doesn't Jesus himself ask, on the cross, ``my God, why have you forsaken me?''

\hfill

\texttt{Albion on 2011-03-10 at 20:04 said:}

Thank you for the response. The question was very poorly phrased indeed; I was hoping for some pointers on where Guenon thought the West had gone off the rails (chronologically), and I see you've posted a link -- it just occured to me that all of the revolutions in the West we've glorified (from 1000 or so onwards) are probably suspect as covert ``aspirations'' which are really a betrayal of unity, and therefore, the Good.

As for Evola, I am a neophyte at all this, but clearly if he had thought himself God, then he would not have felt obligated to clearly condemn so many movements, trends, or fads (eg., homosexuality). It's clear that a fair reading of Evola would support a more positive interpretation, which this website appears to be attempting to do.

I support it.

\hfill

\end{sffamily}
\end{footnotesize}

